\documentclass[10pt,a4paper]{amsart}
\usepackage[margin=19mm]{geometry}
\usepackage{amsmath,amssymb,mathtools}
\usepackage[scr=rsfs,cal=euler]{mathalfa}
\usepackage{xfrac}
\usepackage[unicode,hidelinks,bookmarks=false]{hyperref}
\newcommand{\ie}{i.e.\ }
\newcommand{\cf}{cf.\ }
\newcommand{\resp}{respectively\ }
\newcommand{\Subsection}{\subsection}
\providecommand{\nobreakdash}{\nobreak\hbox{-}}
\setlength{\emergencystretch}{2em}
\multlinegap0pt
\usepackage{amssymb}
\newtheorem{lemma}{Lemma}
\title{Sampling biased spanning forests through marked vertices}
\author{Yves Le Jan}
\date{Source: arXiv:2504.21098v6 (CC BY 4.0)}
\begin{document}
\maketitle

\noindent\emph{Source and licence.} Sampling biased spanning forests through marked vertices. Version arXiv:2504.21098v6 (CC BY 4.0), distributed by arXiv under the Creative Commons Attribution 4.0 International licence, http://creativecommons.org/licenses/by/4.0/, as recorded in the arXiv licence metadata for this exact version and shown on the abstract page https://arxiv.org/abs/2504.21098v6. Full licence notice: \texttt{licenses/arxiv-2504.21098-CC-BY-4.0.txt}.\par\smallskip

\noindent\textit{Editorial scope.} Counted unit: the article's lemma lemdet2 (source0.tex lines 275--277). The article's complete original proof of it is delivered with it (source0.tex lines 327--555, one \texttt{proof} environment of 229 lines, which the article places away from the statement). No mathematics is written, repaired or re-derived: the statement and the proof are the article's own lines, in the article's own notation, and no retained line is substituted or re-flowed.  No further source-local context is needed: every cross-reference of every retained line resolves inside the two delivered spans.  Nothing was omitted from any retained span, so the package carries no omission record.  No retained line contains a citation, so the wrapper carries no bibliography and no bibliography entry.  No retained line contains a figure environment, an image-inclusion command or drawing code, so no image is delivered and the package declares no asset dependency.  Editorial formatting only, in addition: an AMS \texttt{amsart} preamble that restates the article's own declarations and macros used by the retained lines, namely the environments \texttt{lemma} on separate counters, together with the standard packages those lines need.  The retained environments print in this document as Lemma 1 in order of appearance, whereas the article prints the same items with the numbering of its own full document.  The complete native source of the article as distributed in the arXiv e-print for this exact version is bundled unchanged as \texttt{sources/177304/source0.tex}.
 \begin{lemma}\label{lemdet2}
$\mathbb{P}^{(N),\kappa}(\Upsilon_L=Y)=\frac {\kappa^{r-1}(d+\kappa)}{(N+\kappa)^d} $
\end{lemma}

\begin{proof}
Note that a) follows directly from b) since there are $c_l$ binary trees with $l$ labeled leaves.\\
Then we start with the proof of c). Note that setting $\Sigma=\sum_1^{2l-1}[t_i\sqrt N]$
%+l-1$,
 there are $\prod_{i=0}^{\Sigma-1}(N-l-i)$ possible choices for the internal vertices of $\Upsilon_L$, given that $u_L=[t\sqrt N]$.
From lemma \ref{lemdet2}, we get that $$\mathbb P^{(N),\kappa} (Q_L=\alpha, u_L=[t\sqrt N])=\frac{\Sigma+l+\kappa}{(N+\kappa)^{\Sigma+l}}\prod_{i=0}^{\Sigma -1}(N-l-i)$$ (with the usual convention that an empty product equals 1).
 

We define
\[
L_N:=\log(\frac{\Sigma+l+\kappa}{(N+\kappa)^{\Sigma+l}}\prod_{i=0}^{\Sigma-1}(N-l-i)),\]

\begin{lemma}
Set $\Sigma=\sigma\sqrt N$, and assume that
$0\le \sigma\le \sigma_{\max}$, with \(l\) and \(\kappa\) fixed.  
There exists a constant \(C>0\)  (depending on
\(\sigma_{\max},l,\kappa\)) such that for all 
\(\sigma\in[0,\sigma_{\max}]\), and $N$ large enough:
\[
\left|\,L_N-\Bigl(\log(\Sigma+l+\kappa)-l\log N-\frac{\sigma^2}{2}\Bigr)\right|
\le \frac{C}{\sqrt N}.
\]
\end{lemma}
\begin{proof}

\[
L_N=\log(\Sigma+l+\kappa)-(\Sigma+l)\log(N+\kappa)
+\sum_{i=0}^{\Sigma-1}\log(N-l-i),
\]
with $\Sigma=\sum_{i=1}^{2l-1}k_i$ and
$\Sigma\le \sigma_{\max}\sqrt N$.

\medskip
\noindent
Rewrite
\[
\log(N+\kappa)=\log N+\log\!\left(1+\frac{\kappa}{N}\right),
\qquad
\log(N-l-i)=\log N+\log\!\left(1-\frac{l+i}{N}\right),
\]
so that
\[
L_N=\log(\Sigma+l+\kappa)-l\log N
-(\Sigma+l)\log\!\left(1+\frac{\kappa}{N}\right)
+\sum_{i=0}^{\Sigma-1}\log\!\left(1-\frac{l+i}{N}\right).
\]
First

\[
\left|(\Sigma+l)\log\!\left(1+\frac{\kappa}{N}\right)\right|
\le
(\Sigma+l)\frac{\kappa}{N}\le
\frac{\kappa(\sigma_{\max}+l)}{\sqrt N}
:=
\frac{C_1}{\sqrt N}.
\]
%where $C_1=\kappa(\sigma_{\max}+l))$.\\

\medskip
\noindent
%For $0\le x\le 1/2$,
Note that on $(0, 1/2]$,
\[
\log(1-x)=-x-\varepsilon(x)
\quad \text{with}\quad
0\le \varepsilon(x) \le x^2
\]
and set $x_i=(l+i)/N$. Assuming $N$ large enough so that $l+\sigma_{max}\sqrt N\le N/2$,
\[
\sum_{i=0}^{\Sigma-1}\log\!\left(1-\frac{l+i}{N}\right)
=
-\frac1N\sum_{i=0}^{\Sigma-1}(l+i)
-\sum_{i=0}^{\Sigma-1}\varepsilon(x_i).
\]
\medskip
\noindent
We have
\[
0\le \sum_{i=0}^{\Sigma-1}\varepsilon(x_i)
\le
\sum_{i=0}^{\Sigma-1}\left(\frac{l+i}{N}\right)^2
\le
\frac{1}{ N^2}\,\Sigma\,(l+\Sigma)^2
\le
\frac{C_2}{\sqrt N},
\]
with $C_2=\sigma_{\max}(l+\sigma_{\max})^2$.\\

\noindent
Moreover
\[
\frac1N\sum_{i=0}^{\Sigma-1}(l+i)
=
\frac1N\left[\Sigma l+\frac{\Sigma(\Sigma-1)}2\right]=
\frac{\Sigma^2}{2N}
+\frac{\Sigma l}{N}
-\frac{\Sigma}{2N}.
\]
\medskip
\noindent
Hence
\[
\left|
\frac1N\sum_{i=0}^{\Sigma-1}(l+i)
-
\left(\frac{\Sigma^2}{2N}\right)
\right|
\leq
\frac{\Sigma}{2N}+\frac{\Sigma l}{N}
\le
\frac{\sigma_{\max}+2l}{2\sqrt N}
=: \frac{C_3}{\sqrt N}.
\]

\medskip
\noindent
Therefore defining
\[
r_N:=L_N-
\left(\log(\Sigma+l+\kappa)-l\log N-\frac{\Sigma^2}{2N}\right)\]

\[ \left| r_N \right|
\le
\frac{C}{\sqrt N}
\]
with $C=C_1+C_2+C_3$,
uniformly for $\Sigma\le\sigma_{\max}\sqrt N$ and $N$ large enough.
\end{proof}

\medskip
\noindent
%Write
%\[
%\delta_N:=L_N-\Bigl(\log(\Sigma+l+\kappa)-l\log N-\frac{\Sigma^2}{2N}\Bigr),
%\qquad |\delta_N|\le \frac{C}{\sqrt N}.
%\]
Then
\[
e^{L_N}=(\Sigma+l+\kappa)\,N^{-l}e^{-\Sigma^2/(2N)}\,e^{r_N}.
\]
For $N$ large enough so that $C/\sqrt N\le 1$, we have $|e^{r_N}-1|\le 2|r_N|$, hence
\[
\left|e^{L_N}-(\Sigma+l+\kappa)\,N^{-l}e^{-\Sigma^2/(2N)}\right|
\le \frac{2C}{\sqrt N}\,(\Sigma+l+\kappa)\,N^{-l}e^{-\Sigma^2/(2N)},
\]
uniformly for all $\sigma\in[0,\sigma_{\max}]$.
 Moreover, 
 for all $N$ sufficiently large and all admissible
$k\in\mathbb N^{2l-1}$ satisfying $\sum k_i\le \sigma_{\max}\sqrt N$ ( and $k_i >0$ on internal nodes),
\[
\left|
N^{(2l-1)/2}
\mathbb P^{(N),\kappa}(Q_L=\alpha,u_L=k)
-
\left(\sum_{i=1}^{2l-1}\frac{k_i}{\sqrt N}\right)
\exp\!\left(-\frac{(\sum_{i=1}^{2l-1}k_i)^2}{2N}\right)
\right|
\le \frac{C'}{\sqrt N}.
\]
with $C'\le (2C\sigma_{max}+1)$ This concludes the proof of c).\\
\medskip
\noindent
To prove b), we first compute the normalizing integral:
\[
\int_{\mathbb R_+^{2l-1}}
\left(\sum_{i=1}^{2l-1} t_i\right)
\exp\!\left(-\frac{(\sum_{i=1}^{2l-1} t_i)^2}{2}\right)
dt
=
c_l^{-1}.
\]
\medskip
\noindent
Fix $\sigma_{\max}>0$. By the uniform estimate obtained in part (c),
\begin{align*}
\liminf_{N\to\infty}
\mathbb P^{(N),\kappa}(Q_L=\alpha)
&\ge
\liminf_{N\to\infty}
\sum_{\substack{k\in\mathbb N^{2l-1}\\ \sum k_i\le \sigma_{\max}\sqrt N}}
\mathbb P^{(N),\kappa}(Q_L=\alpha,u_L=k)
\\
&=
\liminf_{N\to\infty}
N^{-(2l-1)/2}
\sum_{\substack{k\in\mathbb N^{2l-1}\\ \sum k_i\le \sigma_{\max}\sqrt N}}
\left(\sum_{i=1}^{2l-1}\frac{k_i}{\sqrt N}\right)
\exp\!\left(-\frac{(\sum_{i=1}^{2l-1}k_i)^2}{2N}\right).
\end{align*}
\medskip
\noindent
The last sum is the Riemann sum, with mesh size \(h_N=1/\sqrt N\), over the
simplex
\[
\Delta_{\sigma_{\max}}
:=
\left\{
t\in\mathbb R_+^{2l-1}:\ \sum_{i=1}^{2l-1}t_i\le \sigma_{\max}
\right\}.
\]
Therefore,
\[
\liminf_{N\to\infty}
\mathbb P^{(N),\kappa}(Q_L=\alpha)
\ge
\int_{\Delta_{\sigma_{\max}}} f(t)\,dt .
\]
\medskip
\noindent
Letting $\sigma_{\max}\to\infty$, monotone convergence gives
\[
\liminf_{N\to\infty}
\mathbb P^{(N),\kappa}(Q_L=\alpha)
\ge
\int_{\mathbb R_+^{2l-1}} f(t)\,dt
=
c_l^{-1}.
\]
\medskip
\noindent
Since there are exactly $c_l$ binary shapes, we must have equality:
\[
\lim_{N\to\infty}
\mathbb P^{(N),\kappa}(Q_L=\alpha)
=
c_l^{-1}.
\]
This completes the proof of b).
\end{proof}

\end{document}
